% preamble.tex
% -----------------------------------------------------------------------------
% Single source of truth for packages, theorem environments, macros, and the
% bibliography resource for the KLS frontier manuscript. Inputted by main.tex via
%   \usepackage{subfiles}
%   \input{\subfix{preamble.tex}}
% so that both `latexmk main.tex` and a standalone `latexmk modules/kls/NN-*.tex`
% resolve the same preamble and bibliography.
%
% "Part" in this file always means a \klspart divider of THIS document (Part I ---
% the problem and its frontier, ... , Part V --- synthesis, then the appendices).
% An older revision used the word for the three halves of a since-split unified
% document; no trace of that numbering survives here.
% Do NOT put \documentclass or \usepackage{subfiles} here (those live in
% main.tex).
%
% Build (biber backend), from the repository root:
%   latexmk -pdf -outdir=build main.tex          % the whole document
% A standalone module shows cross-module \ref as "??" and unresolved citations;
% this is expected and harmless for syntax/math checking.
% -----------------------------------------------------------------------------

\usepackage[margin=1in]{geometry}
\usepackage[T1]{fontenc}
\usepackage[utf8]{inputenc}
\usepackage{lmodern}
\usepackage{microtype}
\usepackage{amsmath,amssymb,amsthm,mathtools}
\usepackage{mathrsfs}
\usepackage{bm}
\usepackage{enumitem}
\usepackage{booktabs}
\usepackage{tabularx}
\usepackage{array}
\usepackage{xcolor}
\usepackage{tikz}
\usetikzlibrary{arrows.meta,positioning,fit,backgrounds}
\usepackage{hyperref}
\usepackage[backend=biber,style=alphabetic,maxbibnames=8]{biblatex}
\addbibresource{\subfix{references.bib}}

\hypersetup{
  bookmarksdepth=3,
  colorlinks=true,
  linkcolor=blue!55!black,
  citecolor=blue!55!black,
  urlcolor=blue!55!black
}

% --- claim environments (shared numbering by section) ------------------------
% These ten are the *claim environments*: a \label inside one of them, in modules/,
% must be a ledger node whose `kind` is the environment's name. The environment name
% and the ledger `kind` vocabulary are deliberately the same list, and scripts/check.py
% enforces the correspondence. `remark` is not on the list — it stays expository, which
% is why an obstruction has an environment of its own rather than borrowing one.
\newtheorem{theorem}{Theorem}[section]
\newtheorem{lemma}[theorem]{Lemma}
\newtheorem{proposition}[theorem]{Proposition}
\newtheorem{corollary}[theorem]{Corollary}
\newtheorem{conjecture}[theorem]{Conjecture}
\newtheorem{question}[theorem]{Question}
\theoremstyle{definition}
\newtheorem{definition}[theorem]{Definition}
\newtheorem{assumption}[theorem]{Assumption}
\newtheorem{obstruction}[theorem]{Obstruction}
\theoremstyle{remark}
\newtheorem{example}[theorem]{Example}
\newtheorem{remark}[theorem]{Remark}

% --- program environments (deliberately NOT claim environments) --------------
% A \label inside one of these is structural: scripts/checks/ledger.py resolves an
% anchor's owner against CLAIM_ENVIRONMENTS only, so these need no ledger node and no
% node may claim them. They share the `theorem` counter exactly as the claims do.
% The explicit \theoremstyle{plain} is load-bearing: the block above ends in `remark`
% style, so without the reset `claim` would silently change appearance.
\theoremstyle{plain}
\newtheorem{claim}[theorem]{Claim}
\theoremstyle{definition}
\newtheorem{hypothesis}[theorem]{Hypothesis}
\newtheorem{program}[theorem]{Program}
\newtheorem{family}[theorem]{Family}
\theoremstyle{remark}
\newtheorem{warning}[theorem]{Warning}
\newtheorem{heuristic}[theorem]{Heuristic}

\newcolumntype{Y}{>{\raggedright\arraybackslash}X}

% --- Part dividers -----------------------------------------------------------
% The document is long and multi-route. \klspart{<heading>}{<subtitle>} emits a
% visible divider plus a table-of-contents entry ABOVE \section level, so that the
% grouping of sections into problem / landscape / routes / synthesis / appendices
% is legible both in the TOC and on the page. It introduces no numbering of its
% own: section numbers, theorem numbers, and every \label are untouched, which
% matters because a KLS \label IS a ledger node id (see CLAUDE.md). The Roman
% numeral is therefore part of the heading TEXT, not a counter.
\makeatletter
\newcommand{\l@klspart}[2]{%
  \ifnum \c@tocdepth >-2\relax
    \addpenalty{-\@highpenalty}%
    \vskip 1.0em \@plus\p@
    \begingroup
      \parindent \z@ \rightskip \@pnumwidth
      \parfillskip -\@pnumwidth
      {\leavevmode\bfseries #1}\hfil \nobreak\hb@xt@\@pnumwidth{\hss #2}\par
      \nobreak
    \endgroup
  \fi}
\makeatother
\newcommand{\klspart}[2]{%
  \par\addvspace{2.0\baselineskip}%
  \phantomsection
  \noindent\textcolor{blue!55!black}{\rule{\textwidth}{1.1pt}}\par\vspace{0.30em}%
  {\noindent\large\bfseries #1\par}%
  \vspace{0.15em}%
  {\noindent\itshape #2\par}%
  \vspace{0.30em}%
  \noindent\textcolor{blue!55!black}{\rule{\textwidth}{0.4pt}}\par
  \addcontentsline{toc}{klspart}{#1}%
  \addvspace{1.0\baselineskip}%
}

% --- reader-facing epistemic standing ----------------------------------------
% \klsstatus{<node-id>} sets the standing of that ledger node beside its statement:
% "Published result", "Unreviewed preprint", "Certified here", "Open premise", and so
% on. The strings are NOT written here. They live in status.tex, which is generated from
% research/program/ledger.yaml by `python3 scripts/new.py status`, and
% `check.py --lane editorial` fails the tree if the two disagree --- so the badge on the
% page can never outrun the evidence behind it, and no one has to remember to update it.
%
% The node id is passed explicitly rather than read off the enclosing \label, because
% LaTeX cannot do the latter reliably. The same lane therefore also checks that every
% claim environment carries exactly one \klsstatus naming its own \label.
\input{\subfix{status.tex}}
% Rendered as a small right-flushed tag on its own line directly under the theorem
% header, pulled tight against the statement it belongs to. Colour is decorative: the
% word itself carries the meaning, so the tag survives a greyscale print.
\newcommand{\klsstatus}[1]{%
  \ifcsname klsstatus@#1\endcsname
    \leavevmode\unskip\nobreak\hfill
    {\normalfont\footnotesize\sffamily\color{blue!35!black}%
     [\,\csname klsstatus@#1\endcsname\,]}%
    \par\nobreak\noindent\ignorespaces
  \else
    \PackageWarning{kls}{no recorded standing for '#1'; run scripts/new.py status}%
  \fi}

% --- common macros -----------------------------------------------------------
\newcommand{\R}{\mathbb{R}}
\newcommand{\E}{\mathbb{E}}
\newcommand{\Prob}{\mathbb{P}}
\newcommand{\F}{\mathcal{F}}          % sigma-algebra / filtration
\newcommand{\one}{\mathbf{1}}
\newcommand{\Cov}{\operatorname{Cov}}
\newcommand{\Var}{\operatorname{Var}}
\newcommand{\Tr}{\operatorname{Tr}}
\newcommand{\Ent}{\operatorname{Ent}}
\newcommand{\KL}{\operatorname{KL}}
\newcommand{\Id}{\mathrm{Id}}
\newcommand{\op}{\mathrm{op}}
\newcommand{\osc}{\operatorname{osc}}
\newcommand{\diag}{\operatorname{diag}}
\newcommand{\Hess}{\nabla^2}
\newcommand{\dd}{\,\mathrm{d}}
\newcommand{\eps}{\varepsilon}
\newcommand{\lmax}{\lambda_{\max}}
\newcommand{\lmin}{\lambda_{\min}}

% --- general functional-inequality macros ------------------------------------
\newcommand{\Pstar}{P}                 % target measure P propto e^{-U}
\newcommand{\CP}{C_{\mathrm P}}        % Poincare constant
\newcommand{\CLS}{C_{\mathrm{LS}}}     % log-Sobolev constant
\newcommand{\CTCI}{C_{\mathrm{TCI}}}   % transportation-cost constant
\newcommand{\Wtwo}{W_2}
\newcommand{\Sigmaz}{\Sigma_0}         % Gaussian prior covariance

% --- KLS-specific macros -----------------------------------------------------
\newcommand{\Vol}{\operatorname{Vol}}
\newcommand{\HS}{\mathrm{HS}}
\newcommand{\supp}{\operatorname{supp}}
\newcommand{\dist}{\operatorname{dist}}
\newcommand{\Proj}{\operatorname{Proj}}
\newcommand{\Lip}{\operatorname{Lip}}
\newcommand{\Per}{\operatorname{Per}}
\newcommand{\PsiKLS}{\Psi}
\newcommand{\hstar}{h^{\star}}
\newcommand{\calS}{\mathcal S}
\newcommand{\calT}{\mathcal T}
\newcommand{\calD}{\mathcal D}
\newcommand{\calL}{\mathcal L}
\newcommand{\calI}{\mathcal I}
\newcommand{\calR}{\mathcal R}
\newcommand{\calB}{\mathcal B}
\newcommand{\calQ}{\mathcal Q}
\newcommand{\calK}{\mathcal K}
\newcommand{\calE}{\mathcal E}
\newcommand{\calJ}{\mathcal J}
\newcommand{\calM}{\mathcal M}
\newcommand{\WH}{\mathsf H}
\newcommand{\qJac}{\mathsf q}
\newcommand{\Jsharp}{\calJ^{\sharp}}

% Route C normalization layer (modules/kls/41-*, 42-*): the canonical
% moment-Hessian (CMH) estimate, the affine Poincare constant it dominates,
% and the Stein generator / Dirichlet objects its exact cases use.
\newcommand{\CMH}{C_{\mathrm{CMH}}}       % canonical moment-Hessian constant
\newcommand{\CPaff}{C_{\mathrm P}^{\mathrm{aff}}}  % affine Poincare constant
\newcommand{\Aop}{\mathsf A}              % A = -L_mu, the nonnegative Stein generator
\newcommand{\Div}{\operatorname{div}}
\newcommand{\Dom}{\operatorname{Dom}}
\newcommand{\Dir}{\operatorname{Dir}}
\newcommand{\GammaLaw}{\operatorname{Gamma}}

% --- delimiters --------------------------------------------------------------
\newcommand{\norm}[1]{\left\lVert #1\right\rVert}
\newcommand{\abs}[1]{\left\lvert #1\right\rvert}
\newcommand{\inner}[2]{\left\langle #1,#2\right\rangle}
